\section{Ordered face matching and convex assembly}
\label{cg:assembly}

\subsection{The whole quadrilateral, including its shear}
For cyclically ordered points on a circle or line define
\begin{equation}\label{cg:kappa}
 \kap(a,b,c,d)=\frac{(a-b)(c-d)}{(a-d)(b-c)}>0,
\end{equation}
with limits at infinity. Sending three ordered vertices to
$(\infty,0,1)$ shows that $\kap$ determines the ordered ideal
quadrilateral up to isometry. It fixes its diagonal shear as well
as its labels. A one-step cyclic shift, or reversal fixing the
first vertex, inverts $\kap$. Full reversal
$(a,b,c,d)\mapsto(d,c,b,a)$ preserves it.

The prism coordinates give
\begin{align}
 \kap(F_x)&=\frac{\sin\pi x\sin\pi r}{\sin\pi w\sin\pi p}
              =e^{-m_x},&
 \kap(F_y)&=\frac{\sin\pi y\sin\pi r}{\sin\pi w\sin\pi q}
              =e^{-m_y},\label{cg:input-kappa}\\
 \kap(F_z)&=\frac{\sin\pi p\sin\pi q}{\sin\pi z\sin\pi r}
              =e^{-m_z}.&&\label{cg:output-kappa}
\end{align}
Differentiating \eqref{cg:g} gives the leaf formula
$\kap(C_0,C_1,C_2,C_3)=\tan^2(\pi a)=e^{-g'(a)}$.
For a product $FG$ the actual attachment maps are
\begin{equation}\label{cg:stack-maps}
\begin{array}{c|cccc}
 &C_0&C_1&C_2&C_3\\ \hline
 F&v_0&v_1&P&\infty\\
 G&\infty&P&v_2&v_3 .
\end{array}
\end{equation}
The first is a cyclic shift of $F_x$, and the second is reversal
fixing the first vertex of $F_y$. They invert the input value of
$\kap$. If $f_F$ is the stationary output profile of the lower
child, its wire equation is $f_F'(x)+m_x=0$; the upper equation
is the same with $G,y$. Thus the output and input values of
$\kap$ multiply to one, exactly the condition for the assignments
\eqref{cg:stack-maps} to extend to whole-face isometries.
The new profile has derivative $m_z$. This inductively proves
the output formula $\kap=e^{-f_E'}$.
A turn cyclically relabels the face and replaces $b$ by $1-b$,
which has the same reciprocal effect.

\subsection{Why successive attachments stay convex}
Local angle equations do not by themselves exclude overlap of
distant pieces. The following elementary argument controls the
whole union.

\begin{lemma}\label{cg:convex-union}
Place two finite ideal convex polyhedra on opposite sides of the
plane of an isometric maximal face. If every finite boundary seam
of their union has total interior angle at most $\pi$, the union
is convex. If the common face is a quadrilateral with strictly
nonflat boundary edges in each piece, their ideal closures have
no intersection other than the four recorded face vertices.
\end{lemma}
\begin{proof}
The finite part of the union is closed in hyperbolic space and
locally convex. At a seam edge this is exactly the angle condition;
away from a seam it follows from convexity of the pieces.
Any two finite points can be joined in the union by a finite-length
path. A minimizing sequence with bounded length stays in a compact
ambient hyperbolic ball. Constant-speed parametrization and a
uniformly convergent subsequence give a shortest path in the
union. Local convexity makes it an ambient geodesic near every
point. It is therefore the unique ambient segment, proving
convexity.

The opposite closed half-spaces meet only in their supporting plane.
Their ideal closures can meet only on its boundary circle.
The intersection of either polyhedron with that plane is exactly
its maximal face. An additional ideal vertex in the plane would
enlarge the face across a boundary edge and make that edge flat,
contrary to the strict hypothesis. In the Klein model the closure
of an ideal quadrilateral meets the sphere at infinity at precisely
its four vertices. This excludes every unrecorded ideal
intersection.
\end{proof}

An elementary output has free angle $\pi a$ and cap angle
$\pi(1/2-a)$. Inductively, a completed output $b$ has free angle
$\pi b$ and cap angle $\pi(1-b)$. Attach the lower child first.
At a cap seam a leaf gives total angle $\pi/2$, while a completed
child gives $\pi$ and joins two white faces. The edge shared by the
two prism inputs now has angle
$\pi(w+x)=\pi(1-y)<\pi$. Thus the other input remains a maximal
supported quadrilateral face. Lemma~\ref{cg:convex-union} applies
again when the upper child is attached. We have not assumed
nonintersection between two independently placed children.

The final common input edge has angle $\pi(w+x+y)=\pi$.
The free output seams have angles $\pi(p+x)=\pi z$ and
$\pi(q+y)=\pi z$; the cap output edges have angle $\pi(1-z)$.
Every remaining edge is an original shaded-triangle edge and is
right angled. Flat seams join white faces. Strictness of the
output parameter keeps its face maximal and supported, so the
induction continues.

There is a separate combinatorial check at ideal vertices.
At every open stage the four exposed corners are four unpaired
original terminal points. Product consumes the disjoint pairs
$F_2/G_1$ and $F_3/G_0$ and never exposes them again; a turn merely
permutes the unpaired points. Later identifications cannot extend
a consumed pair into a longer chain. Each final main vertex is
exactly a pair of original terminals, and a leaf's four corners
stay distinct. Its crossing-circle apex is not among those
terminals. Together with Lemma~\ref{cg:convex-union}, this shows
that every leaf retains two distinct shaded triangles sharing only
their original apex.

\subsection{Closing the two completed children}
At the root the outputs are $b,1-b$. The free root equation is
$f_F'(b)=f_G'(1-b)$, so their ordered quadrilaterals have the same
$\kap$. The closing map is
\begin{equation}\label{cg:root}
 C_j^F\longmapsto C_{3-j}^G,\qquad j\pmod4.
\end{equation}
Full reversal preserves $\kap$ and simultaneously realizes the
two stack connections and the two outside closing connections.
Each pair of free angles sums to $\pi$, as does each pair of cap
angles. The last gluing thus joins white faces across flat seams.
The result is a convex right-angled ideal polyhedron $P_E$, with
\begin{equation}\label{cg:half-volume}
 \Vol(P_E)=\pi\sum_\ell g(a_\ell)
             +\pi\sum_{\text{products}}m(x,y,z)=\pi F_0 .
\end{equation}
There are $2H$ shaded triangles, $3H$ ideal vertices and $6H$
edges; Euler's formula gives $H+2$ white faces. These counts are
consequences of the assembly. The next section identifies its
entire labeled boundary.

\section{Recovering the original cut balls}
\label{cg:topology}

Place the pure main strands in the projection sphere, retaining
the crossing disks in their disjoint elementary boxes. Cut the
link exterior along that sphere. Each auxiliary circle becomes
a semicircle in each projection ball, boundary-parallel along its
original half crossing disk. Drilling a boundary-parallel semicircle
introduces one handle; cutting along its half disk opens precisely
that handle. The half disks are disjoint, so these cancellations
can be performed successively in the original boxes. Each result
is a ball with ideal vertices. Main strand segments lie on its
boundary; their removal creates ideal corners and no interior
handles.

At a leaf the cut produces exactly
$(O,C_1,C_2)$ and $(O,C_3,C_0)$ as its shaded triangles.
The other faces are regions of the original projection sphere.
We determine their full walks as follows. At $O$ keep the cyclic
order $(C_0,C_1,C_2,C_3)$. At an outer corner each shaded triangle
contributes the sector ordered by its next and previous neighbors.
The identity strip joining two terminals concatenates their two
sectors in planar order. A boundary walk entering a vertex along
an oriented edge $(u,v)$ exits along the neighbor preceding $u$
in the cyclic order at $v$. Thus the shaded triangles and terminal
pairings determine every remaining white walk.

The maps \eqref{cg:stack-maps} identify $F_2$ with $G_1$ at $P$,
and $F_3$ with $G_0$ at $\infty$. Their exposed order is
$(F_0,F_1,G_2,G_3)=(v_0,v_1,v_2,v_3)$. The prism's two white
triangles insert exactly the intervening white sectors.
At a completed cap seam those sectors merge across a flat
white edge; at a leaf cap seam their boundary remains an original
shaded edge. Turns transport the cyclic order, and
\eqref{cg:root} realizes the original closing pairs. Hence the
ribbon orders of the final nonflat graph are the orders of the
literal projection cut. All faces, including every white walk,
have the same labels and cyclic orders.

This gives a label-preserving cellular homeomorphism of each
cut-ball boundary with $\partial P_E$. A sphere homeomorphism
extends over a ball, for example radially after choosing ball
coordinates. Gluing the two extended balls along their prescribed
white and shaded faces reverses exactly the original cuts.
The quotient is $Y^0_E$, with degree one and the original disk
curves retained.

An ordered quadrilateral during assembly is auxiliary cut data.
It is not the physical four-punctured boundary of the tangle box.
We first used its modulus to assemble cells, and then used the
closed ribbon identification to recover the physical complement.
No boundary group representation is inferred from quadrilateral
moduli alone.

For a retained half twist, cut the twice-punctured disk and rotate
one side through $\pi$. This interchanges both punctures and the
upper and lower half disks. The two balls and white faces are
unchanged. Superscripts distinguishing the projection copies give
the new shaded maps
\begin{equation}\label{cg:odd-pairing}
\begin{array}{c|ccc}
 S_0^+\longrightarrow S_1^-&O&C_1&C_2\\
                          &O&C_3&C_0\\ \hline
 S_1^+\longrightarrow S_0^-&O&C_3&C_0\\
                          &O&C_1&C_2 .
\end{array}
\end{equation}
The other maps are their inverses. Both signs give this unmarked
pairing; the original meridian paths differ, as computed in
Section~\ref{cg:peripheral}. The same labeled balls therefore give
the original $Y^\eps_E$.
